% Standalone theorem insert for W33 papers.
% File: paper/w33_dual_quadrangle_tight_frame_insert.tex

\section{The Dual Quadrangle Tight-Frame and $20$-Fiber Theorem}
\label{sec:dual-quadrangle-tight-frame}

Let $A\in\{-1,0,+1\}^{160\times1620}$ be the signed flag-quadrangle phase matrix.
Rows correspond to the $160$ point-line flags of $W(3,3)$ and columns correspond to
the $1620$ ordinary quadrangles.

\begin{theorem}[Dual Quadrangle Tight-Frame and $20$-Fiber Theorem]
The signed matrix $A$ has rank $81$ and singular spectrum
\[
        A A^T: \quad 160^{81}\oplus0^{79},
        \qquad
        A^T A: \quad 160^{81}\oplus0^{1539}.
\]
Every row has squared norm $81$, and every column has squared norm $8$.
Therefore
\[
        {1\over81}AA^T
\]
is the Gram matrix of $160$ unit flag vectors in an $81$-dimensional phase space,
and
\[
        {1\over8}A^T A
\]
is the Gram matrix of $1620$ unit quadrangle vectors in the same $81$-dimensional
phase space.  The column-frame bound is
\[
        {1620\over81}=20.
\]
Moreover,
\[
        20={1620\over81}={240\over12}={n_B\over h}.
\]
Thus the quadrangle-frame redundancy equals the W33 bulk-to-horizon projection fiber.
\end{theorem}

\begin{proof}[Finite certificate]
The certificate constructs $A$ from the signed $\mathbb F_3$ flag-quadrangle pairing.
It verifies that $\operatorname{rank}(A)=81$, that $AA^T$ has eigenvalue $160$ with
multiplicity $81$ and zero with multiplicity $79$, and that the nonzero spectrum of
$A^TA$ agrees with that of $AA^T$.  Since there are $1620$ quadrangle columns and the
rank is $81$, the normalized column frame has redundancy $1620/81=20$.  The identity
$240/12=20$ identifies this redundancy with the bulk-to-horizon fiber.
\end{proof}

\begin{corollary}[Dual meaning of the signed phase matrix]
The same signed phase matrix simultaneously packages:
\[
\begin{array}{c|c}
\text{side} & \text{tight-frame meaning}\\
\hline
X\text{-flag rows} & 160\text{ local probes spanning }\mathbb R^{81}\\
Z\text{-quadrangle columns} & 1620\text{ loop probes spanning }\mathbb R^{81}\text{ with redundancy }20.
\end{array}
\]
The redundancy $20$ is the holographic fiber $n_B/h$.
\end{corollary}
